%! Populations, Samples, Parameters, and Statistics — Unit 1, Lesson 1.2 — Homework
% PILOT SHAPE (2026-09-06): modeled on AP Statistics lesson 1.4 minus its AP Practice page.
% This is the lesson's SCORED individual practice and the LAST component in the packet.
% Students START IT IN CLASS in the last ten minutes of the period, alone, while the
% teacher circulates for the second formative read; they finish it at home. Packet pages are
% the default; a Desmos activity may replace them on a given day (decided at Close & Assign).
% THIRD CONTEXT: the notes teach on the Lakeside market, the Guided Practice runs on
% Riverbend HS, and these items are on the Harbor Point pool — recall is not the skill.
% Three boxes — items 1-3, 4-5, 6 — so no item breaks across a page (12pt runs three pages).
% Item 4 rows (ii) and (iv) are the deliberate CONTRAST PAIR: two percents, one a statistic
% and one a parameter. Item 6 is the target misconception head-on.
% THE DATA — Harbor Point Community Pool: 1,500 members; 75 asked on a Tuesday evening ->
%   27 lap (36%), 21 family (28%), 18 lessons (24%), 9 aerobics (12%); 0.36 x 1500 = 540;
%   front-desk records: 62% of all 1,500 members renewed; second sample 30/75 = 40%.
% PAGE LOCKSTEP: the key is generated from this file; every work block is byte-identical.
\documentclass[12pt]{article}
\usepackage{saar-article}
\usepackage{saar-boxes}
\newcommand{\TallMath}[1]{$\displaystyle #1\rule[-1.4em]{0pt}{3.2em}$}

\begin{document}

\pageheader{Unit 1, Lesson 1.2}{Homework}
% NAMESTRIP: no \namedateperiod here. The student writes their name once, on the
% cover this component is stapled behind.

\vspace{0.12in}

% Authored identically in homework/ and homework_key/.
\begin{remindbox}
\small
\textbf{This is your graded homework.} It is scored out of 12 and due at the start of the
first class after two study halls. You will start it in the last minutes of the period, so ask about anything that stalls
you before you leave, then finish it on your own.
\end{remindbox}

\vspace{0.08in}

\boxguard[20]
\begin{scenariobox}[Harbor Point Community Pool]{cerulean}
The Harbor Point Community Pool has $1{,}500$ members. The board is deciding whether to
reserve a second lane for lap swimming, so a staff member asks $75$ members arriving on a
Tuesday evening what they mainly use the pool for.

\vspace{2pt}
The $75$ answers: $27$ lap swimming, $21$ family swim, $18$ lessons, $9$ water aerobics.
\end{scenariobox}

\vspace{0.08in}

\boxguard
\begin{notesbox}{Populations, Samples, and the Numbers They Produce}
Every answer names the group it is about --- the $75$ asked, or all $1{,}500$ members.
\begin{enumerate}[leftmargin=1.6em, itemsep=9pt, topsep=3pt]

  \item Name the groups in this study.

        \par\vspace{4pt}
        Population: \blank{5.6cm} \hfill Sample size: \blank{1.6cm}
        \par\vspace{8pt}
        Sample: \blank{6.8cm}

  \item Complete the staff member's frequency table, then show your work for the
        \textbf{water aerobics} row.

        \begin{center}
        \renewcommand{\arraystretch}{1.45}
        \begin{tabularx}{\linewidth}{@{} l c X @{}}
        \toprule
        \textbf{Main use} & \textbf{Count} & \textbf{Relative frequency (percent)} \\
        \midrule
        Lap swimming   & 27 & \blank{6.4cm} \\
        Family swim    & 21 & \blank{6.4cm} \\
        Lessons        & 18 & \blank{6.4cm} \\
        Water aerobics &  9 & \blank{6.4cm} \\
        \midrule
        \textbf{Total} & \blank{1.4cm} & \blank{6.4cm} \\
        \bottomrule
        \end{tabularx}
        \end{center}

        \begin{work}
          \dfrac{9}{75} &= 0.12 \\
                        &= 12\%
        \end{work}

  \item Estimate how many of the pool's $1{,}500$ members mainly swim laps.

        \begin{work}
          0.36 \times 1500 &= 540 \text{ members}
        \end{work}

        Write that result as one sentence for the board. Name the group it is about, and say
        that the number is an estimate.
        \par\writelines{2}

\end{enumerate}
\end{notesbox}

\boxguard[30]
\begin{notesbox}{Populations, Samples, and the Numbers They Produce, continued}
\begin{enumerate}[leftmargin=1.6em, itemsep=8pt, topsep=3pt, start=4]

  \item Label each number \textbf{parameter} or \textbf{statistic}, and say who it describes.

        \begin{center}
        \renewcommand{\arraystretch}{1.5}
        \begin{tabularx}{\linewidth}{@{} X p{2.4cm} >{\raggedright\arraybackslash}p{3.8cm} @{}}
        \toprule
        \textbf{The number} & \textbf{Which kind?} & \textbf{Describes\dots} \\
        \midrule
        The pool has $1{,}500$ members.
          & \blank{2cm} & \blank{3.4cm} \\
        $36\%$ of the $75$ members asked mainly swim laps.
          & \blank{2cm} & \blank{3.4cm} \\
        The true mean number of visits per month for all $1{,}500$ members.
          & \blank{2cm} & \blank{3.4cm} \\
        Front-desk records show $62\%$ of all $1{,}500$ members renewed last year.
          & \blank{2cm} & \blank{3.4cm} \\
        \bottomrule
        \end{tabularx}
        \end{center}

  \item Name the constraint that keeps the board from surveying all $1{,}500$ members in each
        situation.

        \begin{center}
        \renewcommand{\arraystretch}{1.45}
        \begin{tabularx}{\linewidth}{@{} X p{2.8cm} @{}}
        \toprule
        \textbf{Situation} & \textbf{Constraint} \\
        \midrule
        Reaching every member in person would take staff three weeks they do not have.
          & \blank{2.4cm} \\
        Mailing a survey to all $1{,}500$ would cost more than the board budgeted.
          & \blank{2.4cm} \\
        Some members moved away and left no forwarding address.
          & \blank{2.4cm} \\
        Members join and quit every month, so the list is never the same twice.
          & \blank{2.4cm} \\
        \bottomrule
        \end{tabularx}
        \end{center}

\end{enumerate}
\end{notesbox}

\boxguard[30]
\begin{notesbox}{Populations, Samples, and the Numbers They Produce, continued}
\begin{enumerate}[leftmargin=1.6em, itemsep=8pt, topsep=3pt, start=6]

  \item A month later a different staff member asks another $75$ members the same question.
        This time $30$ say they mainly swim laps.

        \textbf{(a)} What percent of that second sample is that?

        \begin{work}
          \dfrac{30}{75} &= 0.40 \\
                         &= 40\%
        \end{work}

        \textbf{(b)} The two surveys gave $36\%$ and $40\%$. A board member says, ``One of you
        counted wrong.'' Explain why the difference does \emph{not} mean that, and say how many
        true percents there are for all $1{,}500$ members.
        \par\writelines{2}

        \textbf{(c)} The board's draft report says: \textit{``Harbor Point's $1{,}500$ members
        want a second lap lane.''} Every count behind it is correct. Explain \emph{two} separate
        things that sentence claims which the survey cannot support.
        \par\writelines{2}

\end{enumerate}
\end{notesbox}

\vspace{0.08in}

\boxguard[10]
\begin{spiralbox}
\textbf{Next lesson: Choosing a Sample --- Four Sampling Techniques.} The staff member asked
whoever showed up on a Tuesday evening --- and lap swimmers keep evening hours. Next class is
about \emph{how} the sample gets chosen, and why that choice decides whether a statistic is worth
anything at all.
\end{spiralbox}

\end{document}
